% GENERATED FILE. Edit canonical lessons, then run npm run generate:books.
% canonical-source-hash: fnv1a64:570f7b5bcadcf145
% canonical-lessons: TA-C396-tuntu, TA-C396-etirvinai, TA-C396-cumattu, TA-C396-piratipali, TA-C396-porutpatuttu

\chapter{Pressure and Response}
\label{ch:ta-pressure-and-response}
\hlchaptermodality{396}

\begin{chapteropening}
\textbf{By the end of this chapter:} \emph{I can talk about stirring people up and how they react.}

Complete the last lesson of chapter 396: I can talk about stirring people up and how they react.
\end{chapteropening}

\section[\texorpdfstring{\ta{தூண்டு} (tūṇṭu)}{தூண்டு (tūṇṭu)}]{\ta{தூண்டு} (tūṇṭu) — to incite; to stir up}

\label{lesson:TA-C396-tuntu}



\begin{quote}
Before the new one: say the Tamil for to abolish, then the Tamil for to abandon.
\end{quote}

\subsection*{You'll want to know: \ta{தூண்டு}}
\textbf{\ta{தூண்டு}} — \emph{tūṇṭu} — \textquotedblleft{}to incite; to stir up\textquotedblright{}.

\begin{quote}
\ta{அவருடைய} \ta{பேச்சு} \ta{மக்களைப்} \ta{போராடத்} \ta{தூண்டியது}.
\end{quote}

\begin{quote}
\emph{avaruṭaiya pēccu makkaḷaip pōrāṭat tūṇṭiyatu.}
\end{quote}

\begin{quote}
His speech stirred people to fight.
\end{quote}

Past (I …): \ta{தூண்டினேன்}, \emph{tūṇṭiṉēṉ}.

Who is moved takes -\ta{ஐ} and the action an infinitive: \ta{படிக்கத்} \ta{தூண்டு}, to spur someone to study.

\begin{cousinweb}
\textbf{\ta{தூண்டு}} (\emph{tūṇṭu}) is native Dravidian: DEDR 3380 lists Tamil \emph{tūṇṭu}, \textquotedblleft{}to shoot, propel; to incite, goad, remind, suggest\textquotedblright{}.
\end{cousinweb}

\begin{grammarlens}[title={What you've built}]
One more word for saying what happens if.
\end{grammarlens}

\subsection*{Your turn}
\begin{itemize}
\raggedright
  \item \emph{Say it:} \emph{tūṇṭu}
  \item \emph{Say it:} \emph{tūṇṭu}, once more
  \item \emph{Say it:} \emph{avaruṭaiya pēccu makkaḷaip pōrāṭat tūṇṭiyatu.} — His speech stirred people to fight.
  \item \emph{Recall:} say the Tamil for to abolish, then the Tamil for to abandon, then say \emph{tūṇṭu} again
\end{itemize}

\begin{tcolorbox}[breakable,colback=teal!4,colframe=teal!35!black,title={Before you move on}]
What does \ta{தூண்டு} mean? (\textquotedblleft{}To incite; to stir up\textquotedblright{}.) Say it once more.
\end{tcolorbox}
\section[\texorpdfstring{\ta{எதிர்வினை} (etirviṉai)}{எதிர்வினை (etirviṉai)}]{\ta{எதிர்வினை} (etirviṉai) — a reaction; a response}

\label{lesson:TA-C396-etirvinai}



\begin{quote}
Before the new one: say the Tamil for to abandon, then the Tamil for to incite.
\end{quote}

\subsection*{You'll want to know: \ta{எதிர்வினை}}
\textbf{\ta{எதிர்வினை}} — \emph{etirviṉai} — \textquotedblleft{}a reaction; a response\textquotedblright{}.

\begin{quote}
\ta{இந்த} \ta{முடிவுக்கு} \ta{மக்களிடமிருந்து} \ta{பெரிய} \ta{எதிர்வினை} \ta{வந்தது}.
\end{quote}

\begin{quote}
\emph{inta muṭivukku makkaḷiṭamiruntu periya etirviṉai vantatu.}
\end{quote}

\begin{quote}
There was a big reaction from the people to this decision.
\end{quote}

\ta{எதிர்}, against + \ta{வினை}, action: action in return.

\begin{cousinweb}
\textbf{\ta{எதிர்வினை}} (\emph{etirviṉai}) is a modern native coinage: \emph{etir}, \textquotedblleft{}opposite, in return\textquotedblright{} (DEDR 795), and \emph{viṉai}, \textquotedblleft{}action, deed\textquotedblright{} (DEDR 5445).
\end{cousinweb}

\begin{grammarlens}[title={What you've built}]
One more word for saying what happens if.
\end{grammarlens}

\subsection*{Your turn}
\begin{itemize}
\raggedright
  \item \emph{Say it:} \emph{etirviṉai}
  \item \emph{Say it:} \emph{etirviṉai}, once more
  \item \emph{Say it:} \emph{inta muṭivukku makkaḷiṭamiruntu periya etirviṉai vantatu.} — There was a big reaction from the people to this decision.
  \item \emph{Recall:} say the Tamil for to abandon, then the Tamil for to incite, then say \emph{etirviṉai} again
\end{itemize}

\begin{tcolorbox}[breakable,colback=teal!4,colframe=teal!35!black,title={Before you move on}]
What does \ta{எதிர்வினை} mean? (\textquotedblleft{}A reaction; a response\textquotedblright{}.) Say it once more.
\end{tcolorbox}
\section[\texorpdfstring{\ta{சுமத்து} (cumattu)}{சுமத்து (cumattu)}]{\ta{சுமத்து} (cumattu) — to impose; to lay on}

\label{lesson:TA-C396-cumattu}



\begin{quote}
Before the new one: say the Tamil for to incite, then the Tamil for a reaction.
\end{quote}

\subsection*{You'll want to know: \ta{சுமத்து}}
\textbf{\ta{சுமத்து}} — \emph{cumattu} — \textquotedblleft{}to impose; to lay on\textquotedblright{}.

\begin{quote}
\ta{அதிக} \ta{வரி} \ta{சுமத்தினால்} \ta{மக்கள்} \ta{கோபப்படுவார்கள்}.
\end{quote}

\begin{quote}
\emph{atika vari cumattiṉāl makkaḷ kōpappaṭuvārkaḷ.}
\end{quote}

\begin{quote}
If heavy taxes are imposed, people will get angry.
\end{quote}

Past (I …): \ta{சுமத்தினேன்}, \emph{cumattiṉēṉ}.

The load stands bare or takes -\ta{ஐ}, the person \ta{மீது}: X \ta{மீது} \ta{பழி} \ta{சுமத்து}, to lay the blame on X.

\begin{cousinweb}
\textbf{\ta{சுமத்து}} (\emph{cumattu}) is native Dravidian: DEDR 2677 lists Tamil \emph{cuma}, \textquotedblleft{}to bear, carry a burden\textquotedblright{}, and its causative \emph{cumattu}, \textquotedblleft{}to load on\textquotedblright{}.
\end{cousinweb}

\begin{grammarlens}[title={What you've built}]
One more word for saying what happens if.
\end{grammarlens}

\subsection*{Your turn}
\begin{itemize}
\raggedright
  \item \emph{Say it:} \emph{cumattu}
  \item \emph{Say it:} \emph{cumattu}, once more
  \item \emph{Say it:} \emph{atika vari cumattiṉāl makkaḷ kōpappaṭuvārkaḷ.} — If heavy taxes are imposed, people will get angry.
  \item \emph{Recall:} say the Tamil for to incite, then the Tamil for a reaction, then say \emph{cumattu} again
\end{itemize}

\begin{tcolorbox}[breakable,colback=teal!4,colframe=teal!35!black,title={Before you move on}]
What does \ta{சுமத்து} mean? (\textquotedblleft{}To impose; to lay on\textquotedblright{}.) Say it once more.
\end{tcolorbox}
\section[\texorpdfstring{\ta{பிரதிபலி} (piratipali)}{பிரதிபலி (piratipali)}]{\ta{பிரதிபலி} (piratipali) — to reflect; to mirror}

\label{lesson:TA-C396-piratipali}



\begin{quote}
Before the new one: say the Tamil for a reaction, then the Tamil for to impose.
\end{quote}

\subsection*{You'll want to know: \ta{பிரதிபலி}}
\textbf{\ta{பிரதிபலி}} — \emph{piratipali} — \textquotedblleft{}to reflect; to mirror\textquotedblright{}.

\begin{quote}
\ta{இந்த} \ta{முடிவு} \ta{மக்களின்} \ta{கருத்தைப்} \ta{பிரதிபலிக்கிறது}.
\end{quote}

\begin{quote}
\emph{inta muṭivu makkaḷiṉ karuttaip piratipalikkiṟatu.}
\end{quote}

\begin{quote}
This decision reflects the view of the people.
\end{quote}

Past (it …): \ta{பிரதிபலித்தது}, \emph{piratipalittatu}.

Also of light: \ta{ஏரி} \ta{நிலாவைப்} \ta{பிரதிபலிக்கிறது}, the lake reflects the moon.

\begin{cousinweb}
\textbf{\ta{பிரதிபலி}} (\emph{piratipali}) is a verb built on Sanskrit \emph{pratiphala}, \textquotedblleft{}reflection, image, shadow\textquotedblright{} (Monier-Williams), with the loan-verb ending -\ta{இ}.
\end{cousinweb}

\begin{grammarlens}[title={What you've built}]
One more word for saying what happens if.
\end{grammarlens}

\subsection*{Your turn}
\begin{itemize}
\raggedright
  \item \emph{Say it:} \emph{piratipali}
  \item \emph{Say it:} \emph{piratipali}, once more
  \item \emph{Say it:} \emph{inta muṭivu makkaḷiṉ karuttaip piratipalikkiṟatu.} — This decision reflects the view of the people.
  \item \emph{Recall:} say the Tamil for a reaction, then the Tamil for to impose, then say \emph{piratipali} again
\end{itemize}

\begin{tcolorbox}[breakable,colback=teal!4,colframe=teal!35!black,title={Before you move on}]
What does \ta{பிரதிபலி} mean? (\textquotedblleft{}To reflect; to mirror\textquotedblright{}.) Say it once more.
\end{tcolorbox}
\section[\texorpdfstring{\ta{பொருட்படுத்து} (poruṭpaṭuttu)}{பொருட்படுத்து (poruṭpaṭuttu)}]{\ta{பொருட்படுத்து} (poruṭpaṭuttu) — to take into account; to mind}

\label{lesson:TA-C396-porutpatuttu}



\begin{quote}
Before the new one: say the Tamil for to impose, then the Tamil for to reflect.
\end{quote}

\subsection*{You'll want to know: \ta{பொருட்படுத்து}}
\textbf{\ta{பொருட்படுத்து}} — \emph{poruṭpaṭuttu} — \textquotedblleft{}to take into account; to mind\textquotedblright{}.

\begin{quote}
\ta{அவன்} \ta{மழையைப்} \ta{பொருட்படுத்தாமல்} \ta{வெளியே} \ta{போனான்}.
\end{quote}

\begin{quote}
\emph{avaṉ maḻaiyaip poruṭpaṭuttāmal veḷiyē pōṉāṉ.}
\end{quote}

\begin{quote}
He went out regardless of the rain.
\end{quote}

Past (I …): \ta{பொருட்படுத்தினேன்}, \emph{poruṭpaṭuttiṉēṉ}.

Mostly negative: regardless of, \ta{பொருட்படுத்தாமல்}; never mind, \ta{பொருட்படுத்த} \ta{வேண்டாம்}.

\begin{cousinweb}
\textbf{\ta{பொருட்படுத்து}} (\emph{poruṭpaṭuttu}) is native: \emph{poruḷ}, \textquotedblleft{}thing, meaning\textquotedblright{} (DEDR 4544), whose ḷ becomes \ta{ட்} before p, with \emph{paṭuttu}: to treat something as a thing that matters.
\end{cousinweb}

\begin{grammarlens}[title={What you've built}]
One more word for saying what happens if.
\end{grammarlens}

\subsection*{Your turn}
\begin{itemize}
\raggedright
  \item \emph{Say it:} \emph{poruṭpaṭuttu}
  \item \emph{Say it:} \emph{poruṭpaṭuttu}, once more
  \item \emph{Say it:} \emph{avaṉ maḻaiyaip poruṭpaṭuttāmal veḷiyē pōṉāṉ.} — He went out regardless of the rain.
  \item \emph{Recall:} say the Tamil for to impose, then the Tamil for to reflect, then say \emph{poruṭpaṭuttu} again
\end{itemize}

\begin{tcolorbox}[breakable,colback=teal!4,colframe=teal!35!black,title={Before you move on}]
What does \ta{பொருட்படுத்து} mean? (\textquotedblleft{}To take into account; to mind\textquotedblright{}.) Say it once more.
\end{tcolorbox}
